\chapter*{Resumo}
\thispagestyle{empty}

O diabetes mellitus, com ênfase crítica no Tipo 2, consolidou-se como uma das emergências sanitárias mais prementes do cenário global, sendo uma doença crônica intrinsecamente associada a fatores de estilo de vida modernos que impõe uma carga severa e crescente aos sistemas de saúde pública. Diante da vastidão e heterogeneidade dos dados gerados no acompanhamento dessa patologia, as abordagens analíticas tradicionais tornam-se muitas vezes limitadas, o que justifica a emergência das redes neurais artificiais como ferramentas tecnológicas indispensáveis e promissoras para a detecção e prevenção. Diferentemente de métodos estatísticos lineares, esses algoritmos avançados possuem a capacidade de processar grandes volumes de informações clínicas, demográficas e comportamentais, identificando padrões complexos, ocultos e não lineares para classificar com alta precisão o risco individual de desenvolvimento da doença. A integração estratégica desses modelos preditivos em sistemas de apoio à decisão médica representa um avanço significativo, pois não apenas potencializa a realização de diagnósticos mais rápidos e assertivos, mas também fundamenta a medicina de precisão através do desenvolvimento de estratégias de intervenção precoce personalizadas. Em última análise, a aplicação dessa inteligência computacional visa reverter as tendências epidemiológicas atuais, reduzindo drasticamente a incidência, as complicações a longo prazo e o impacto socioeconômico do diabetes na sociedade contemporânea.
\\\\
Palavras chave: Redes neurais, diabetes, predição, inteligência artificial, aprendizado de máquina.